% ===== Envoi / Blank =====
\section*{Envoi}
\addcontentsline{toc}{section}{Envoi}
\label{p19:sec:envoi}

\noindent
By \xref{p19:sub:unwrit-close}, the fitting end is not a conclusion. What follows is not a summary, makes no further claim, and states no thesis---any added signifier would weaken the needlessness of \xref{p19:sub:silence-three} and the laying-down of the pen of \xref{p19:sub:unwrit-close}. The envoi performs only an \emph{indication} (a turning-toward, not a statement) and an \emph{unclosing image}. It does not speak \emph{about} happiness; it enacts having no more need to speak. The ``we'' is left anonymous---a bare structure, into which any reader may enter their own afternoon.

\vspace{1.5em}

\noindent
And so we return the gaze to the afternoon that is, even now, still going on. Two people. The rain beyond the window. Nothing to be done, and nothing lacked.

\vspace{1.0em}

% --- The closing image is to be set by the author, by hand. ---
% Discipline (per the paper): one unclosing image, no recapitulation, no
% return to any thesis term ("generativity," "happiness," etc.); meaning
% moved onto the path of reading, extending anew with each re-reading,
% refusing closure (cf. Paper IX). Reticent, image-based, in the author's
% own voice. Left deliberately unset here, to be completed by Wanhong.

\begin{center}
\itshape\color{rose}
[\,the closing image, to be set by hand\,]
\end{center}

\vspace{2.0em}

% --- Blank to end: the typographic instance of the third silence. ---
% The paper stops at the boundary of what can be said, adding no further
% signifier. The blank of the page is the silence of needlessness made
% visible (cf. §7.2). Leave the remainder of the page empty.

\vfill
\noindent\hfill\textcolor{gold}{\rule{0.2\linewidth}{0.4pt}}\hfill\null
